% !PRINTED 566
\weberpage{566}

\begin{center}
\S.\,156.

\gesperrt{Anwendung auf algebraische K\"orper.}
\end{center}

Aus dem Minimum f\"ur eine positive quadratische Form kann
man ein Minimum von Linearformen ableiten, indem man eine
positive quadratische Form als Summe von Quadraten von Linearformen
 auffasst. Darauf beruhen die folgenden Anwendungen
auf die Theorie der algebraischen K\"orper, denen wir einen H\"ulfssatz
 vorausschicken:

1. \gesperrt{Sind $a_1, a_2, \ldots, a_n$ reelle positive Zahlwerthe, so ist}
\weq{1}{a_1 + a_2 + \cdots + a_n \geqq n \sqrt[n]{a_1 a_2 \ldots a_n}.}

Dieser Satz ist offenbar richtig f\"ur $n=2$, denn es ist
$a_1 + a_2 - 2\sqrt{a_1 a_2} = (\sqrt{a_1} - \sqrt{a_2})^2$, also nie negativ. Wir
wenden also die vollst\"andige Induction an, indem wir die Formel
(1) f\"ur $n-1$ Glieder als schon erwiesen annehmen. Nehmen
wir dann an, dass $a_n$ von keinem der \"ubrigen $a$ an Gr\"osse \"ubertroffen
 wird, so ist
\[ a_n \geqq \sqrt[n-1]{a_1 a_2 \ldots a_{n-1}} = b, \]
und
\[ a_1 + a_2 + \cdots + a_{n-1} \geqq (n-1)\,b. \]
Also ist
\[ a_1 + a_2 + \cdots + a_n \geqq a_n + (n-1)\,b, \]
und folglich
\[ \frac{a_1 + a_2 + \cdots + a_n}{n} \geqq b + \frac{a_n - b}{n}, \]

\begin{footnotesize}
\noindent\rule{4cm}{0.4pt}

worin $\Gamma(x)$ die Gammafunction bedeutet, f\"ur die die Gleichungen
\[ \Gamma(x) = (x-1)\,\Gamma(x-1), \quad \Gamma\left(\frac{1}{2}\right) = \sqrt{\pi} \]
bestehen, so dass $K_n$ durch rationale Zahlen und durch $\sqrt{\pi}$ ausgedr\"uckt ist.

(Vergl. z.\,B. Vorlesungen \"uber die Theorie der bestimmten Integrale
nach \gesperrt{Lejeune-Dirichlet} von G.\,F. \gesperrt{Meyer} herausgegeben, \S.\,175.
Leipzig 1871).

Daraus ergiebt sich
\[ \Theta_n = \frac{4}{n\pi}\,\Gamma\left(1 + \frac{n}{2}\right)^{\frac{2}{n}}, \]
was nach den N\"aherungsformeln f\"ur die $\Gamma$-Function f\"ur grosse
Werthe von $n$ nahe mit
\[ \frac{2}{\pi e} = 0{,}234\ldots \]
\"ubereinstimmt. (\gesperrt{Minkowski}, Crelle, Bd. 107, S. 292).
\end{footnotesize}
%CONT>>
